% Standalone section; no preamble.  Standard amsmath/amssymb/amsthm are sufficient.
% Suggested bibliography keys appear below.  This is a self-contained
% specialization of established q-difference resurgence, not a claim that
% resurgence of the q-Pochhammer symbol is newly discovered.

\section{Cyclotomic Borel transforms and the necessary Stokes correction}
\label{sec:cyclotomic-borel-correction}

The repository's conjecture \texttt{conj:cyclotomic-resurgent-inverse} \cite{repo-qfrontier}
contains two different issues: resurgence of the normalized logarithm of
the product, and identification of all its directional sums with one fixed
analytic function.  The first issue admits an explicit answer throughout
the stated disk $|a|<1$.  The second identification is false as written:
different nonsingular Laplace rays can give different functions even where
both integrals converge and the original product is defined.  The
correction is an explicitly computable exponentially small term.

This analysis belongs to an established literature.  Marmi and Sauzin
\cite[\S4.2--4.3, especially Theorem~4.3]{marmi-sauzin} study Borel
transforms at arbitrary root-of-unity resonances for the cohomological
equation.  Their general solution
$f_g(q,z)=\sum_{r\geq1}g_r z^r/(q^r-1)$ includes the present logarithm
when $g(z)=-\log(1-z)$.  Their variable is $\eta=\log(q/\zeta)=-t$,
and their theorem is formulated for $\eta f_g$; differentiation in the
Borel variable gives the normalization used below.  Garoufalidis and
Kashaev \cite{garoufalidis-kashaev} give explicit poles, residues, and
Stokes transformations for Faddeev's quantum dilogarithm.  The more recent
work of Fantini and Rella \cite[\S3]{fantini-rella} studies
$q$-Pochhammer resurgence and the correction terms needed to reconstruct
particular $q$-series.  Our purpose is a directly checkable specialization,
with uniform error bounds and a concrete repair of the repository's
normalization claim.

\subsection{Normalization and the exact Borel transform}

Let $m\geq1$, let $h$ be an integer coprime to $m$, and put
$\zeta=\exp(2\pi i h/m)$.  For $|a|<1$ and $\Re t>0$, define
\begin{equation}
 L_a(t)=\log(a;\zeta e^{-t})_\infty
       =-\sum_{r=1}^{\infty}
          \frac{a^r}{r(1-\zeta^r e^{-rt})}.
 \label{eq:cb-log-definition}
\end{equation}
Here and below the logarithm is defined by its convergent series.  Set
\begin{align}
 S_{-1}(a)&=-\frac{\operatorname{Li}_2(a^m)}{m^2},\label{eq:cb-action}\\
 S_0(a)&=\frac{1}{2m}\log(1-a^m)
       -\sum_{\substack{r\geq1\\m\nmid r}}
          \frac{a^r}{r(1-\zeta^r)},\label{eq:cb-constant}\\
 R_a(t)&=L_a(t)-\frac{S_{-1}(a)}{t}-S_0(a).
 \label{eq:cb-normalization}
\end{align}
For $j\in\mathbb Z\setminus\{0\}$, let
\begin{equation}
 \omega_j=\frac{2\pi i j}{m},\qquad
 p_j\in\{1,\ldots,m\},\qquad hp_j\equiv j\pmod m.
 \label{eq:cb-indices}
\end{equation}
The convention $p_j=m$ represents the zero residue class.

\begin{theorem}[Explicit cyclotomic Borel transform]
\label{thm:cyclotomic-borel}
The normalized function $R_a(t)$ has a formal expansion
$\widetilde R_a(t)=\sum_{n\geq1}c_n(a)t^n$ with
\begin{equation}
 c_n(a)=\sum_{r\geq1}a^r r^{n-1}
       \sum_{\substack{j\ne0\\j\equiv hr\pmod m}}
          \omega_j^{-n-1}\qquad(n\geq1).
 \label{eq:cb-coefficients}
\end{equation}
Its Borel transform, with the convention
$\mathcal B(t^n)=\xi^{n-1}/(n-1)!$, is
\begin{equation}
 B_{a,\zeta}(\xi)
 =\sum_{j\ne0}\frac{1}{\omega_j^2}
   \frac{a^{p_j}\exp(p_j\xi/\omega_j)}
        {1-a^m\exp(m\xi/\omega_j)}.
 \label{eq:cb-borel}
\end{equation}
For $a\ne0$, choose any $\lambda$ with $a=e^{-\lambda}$; necessarily
$\Re\lambda>0$.  The complete set of singularities is
\begin{equation}
 \Sigma_{a,\zeta}
 =\left\{\xi_{j,n}:=
       \omega_j\left(\lambda+\frac{2\pi i n}{m}\right):
       j\in\mathbb Z\setminus\{0\},\ n\in\mathbb Z\right\}.
 \label{eq:cb-poles}
\end{equation}
These points are distinct, nonzero, and locally finite.  Every singularity
is a simple pole, with
\begin{equation}
 \operatorname*{Res}_{\xi=\xi_{j,n}} B_{a,\zeta}(\xi)
 =-\frac{\exp(2\pi i n p_j/m)}{m\omega_j}.
 \label{eq:cb-residues}
\end{equation}
Thus the Borel germ extends as a single-valued meromorphic function on
$\mathbb C$, in particular along every path avoiding
$\Sigma_{a,\zeta}$.  For $a=0$ it is identically zero.
\end{theorem}

\begin{proof}
For a fixed $r\geq1$, let
$J_r=\{j\ne0:j\equiv hr\pmod m\}$ and put
\[
 \kappa_r=\begin{cases}
  1/2,&m\mid r,\\
  (1-\zeta^r)^{-1},&m\nmid r.
 \end{cases}
\]
The Mittag--Leffler expansion of the hyperbolic cotangent, after
subtraction of its value at zero and removal of a possible pole at zero,
gives
\begin{equation}
 \frac{1}{1-\zeta^r e^{-z}}
 =\frac{{\bf1}_{m\mid r}}{z}+\kappa_r
  -z\sum_{j\in J_r}
       \frac{1}{\omega_j^2(1-z/\omega_j)}.
 \label{eq:cb-kernel}
\end{equation}
For clarity, the nonconstant part of the partial-fraction expansion is
$\sum_{j\in J_r}[(z-\omega_j)^{-1}+\omega_j^{-1}]$,
which equals the last term in \eqref{eq:cb-kernel}.  This subtracted series
converges normally on compact sets away from its poles.  The identity
also follows immediately from
$1/(1-e^{-w})=1/2+(1/2)\coth(w/2)$ after a purely imaginary translation.

Substitute $z=rt$ in \eqref{eq:cb-kernel}, multiply by $-a^r/r$, and
sum over $r$.  The pole and constant contributions are precisely
\eqref{eq:cb-action} and \eqref{eq:cb-constant}.  Consequently
\begin{equation}
 R_a(t)=t\sum_{r\geq1}a^r
             \sum_{j\in J_r}
               \frac{1}{\omega_j^2(1-rt/\omega_j)}.
 \label{eq:cb-exact-kernel}
\end{equation}
All relevant series converge absolutely for each $\Re t>0$; this follows,
for example, from the denominator estimate in the next theorem with a
closed subsector containing $t$.  Expanding the rational factor formally
gives \eqref{eq:cb-coefficients}.  Applying the Borel transform gives
\[
 B_{a,\zeta}(\xi)
 =\sum_{r\geq1}a^r\sum_{j\in J_r}
          \omega_j^{-2}e^{r\xi/\omega_j}.
\]
In a sufficiently small disk about zero, these series converge
absolutely.  For a given $j$, the permitted $r$ are $p_j+\ell m$,
$\ell\geq0$; their geometric sum is \eqref{eq:cb-borel}.

The denominator of its $j$th summand vanishes exactly when
$\xi/\omega_j=\lambda+2\pi i n/m$.  At such a point the numerator is
$\exp(2\pi i n p_j/m)$ and the derivative of the denominator is
$-m/\omega_j$, proving \eqref{eq:cb-residues}.  Moreover
\[
 \Im\xi_{j,n}=\frac{2\pi j}{m}\Re\lambda.
\]
In a bounded disk this bounds $|j|$, and then bounds $|n|$.  It also shows
that equality of two locations forces equality first of $j$, then of
$n$.  In particular, there is no cancellation between summands at a pole.
Finally the tail of \eqref{eq:cb-borel} is locally $O(j^{-2})$, since
$\xi/\omega_j\to0$ and $1-a^m\ne0$.  This proves normal convergence away
from the displayed poles and completes the continuation argument.
Changing the logarithm $\lambda$ merely reindexes $n$.
\end{proof}

\begin{corollary}[Exact convergence radius and large-order coefficients]
\label{cor:cb-large-order}
For $a\ne0$, the convergence radius of the Borel germ is exactly
\begin{equation}
 d(a,m)=\frac{2\pi}{m}
        \min_{n\in\mathbb Z}\left|\lambda+\frac{2\pi i n}{m}\right|.
 \label{eq:cb-exact-radius}
\end{equation}
It depends on the order $m$ but not on the choice of primitive $m$th
root $\zeta$.  The formal expansion has optimal Gevrey order one: it
does not have Gevrey order $s<1$.
Let $P_0$ be the finite set of poles of modulus $d(a,m)$, let $d_1>d(a,m)$
be the next pole modulus, and fix $d(a,m)<b<d_1$.  Then
\begin{equation}
 \frac{c_n(a)}{(n-1)!}
 =-\sum_{\sigma\in P_0}
       \frac{\operatorname{Res}_{\xi=\sigma}B_{a,\zeta}(\xi)}{\sigma^n}
   +O\!\left(b^{-(n-1)}\right)
 \qquad(n\to\infty).
 \label{eq:cb-large-order}
\end{equation}
The finite sum is kept intact; it may cancel on subsequences, as it must
for the vanishing even coefficients when $m=1$.
\end{corollary}

\begin{proof}
Minimizing \eqref{eq:cb-poles} first over $j$ gives $|j|=1$, proving
\eqref{eq:cb-exact-radius}.  None of these poles is removable by
\eqref{eq:cb-residues}.  If $|c_n|\leq CA^n(n!)^s$ for some $s<1$,
the Borel transform would have infinite convergence radius, a
contradiction.  Finally subtract from $B_{a,\zeta}$ the finitely many
principal parts $\operatorname{Res}_{\sigma}B/(\xi-\sigma)$ for
$\sigma\in P_0$.  The difference is holomorphic on $|\xi|<d_1$.
Cauchy's coefficient bound on $|\xi|=b$ and the geometric expansion of
each subtracted principal part give \eqref{eq:cb-large-order}.
\end{proof}

\subsection{Uniform Gevrey bounds and reconstruction of the product}

\begin{theorem}[Uniform remainders and positive-ray reconstruction]
\label{thm:cb-gevery-laplace}
Fix $0<\rho<1$ and $0<\delta<\pi/2$, and write
$\ell=-\log\rho$.  Uniformly for $|a|\leq\rho$ and
$|\arg t|\leq\pi/2-\delta$, for every integer $N\geq1$,
\begin{equation}
 \left|R_a(t)-\sum_{n=1}^{N-1}c_n(a)t^n\right|
 \leq
 \frac{m\pi}{6\rho\sin\delta}
 \left(\frac{m}{2\pi\ell}\right)^N
 (N-1)!\,|t|^N.
 \label{eq:cb-gevery-bound}
\end{equation}
The Borel transform is uniformly bounded on the real axis:
\begin{equation}
 \sup_{\xi\in\mathbb R}|B_{a,\zeta}(\xi)|
 \leq \frac{m^2\rho}{12(1-\rho^m)}.
 \label{eq:cb-real-bound}
\end{equation}
More generally it is holomorphic and uniformly bounded on every closed
horizontal strip $|\Im\xi|\leq b$ with
$b<2\pi\ell/m$.  For all $\Re t>0$,
\begin{equation}
 R_a(t)=\int_0^{\infty}e^{-\xi/t}B_{a,\zeta}(\xi)\,d\xi.
 \label{eq:cb-positive-laplace}
\end{equation}
\end{theorem}

\begin{proof}
If $|\arg t|\leq\pi/2-\delta$, then the argument of $rt/\omega_j$
is separated from the positive real axis by at least $\delta$.  Hence
\[
 |1-rt/\omega_j|\geq\sin\delta
 \qquad(r\geq1,\ j\ne0).
\]
Use the finite geometric identity
$ (1-u)^{-1}=\sum_{k=0}^{N-2}u^k+u^{N-1}(1-u)^{-1}$
in \eqref{eq:cb-exact-kernel}; for $N=1$ the finite sum is empty.
The exact remainder equals
\[
 t^N\sum_{r\geq1}a^r r^{N-1}
       \sum_{j\in J_r}
        \frac{\omega_j^{-N-1}}{1-rt/\omega_j}.
\]
Its absolute value is at most
\[
 \frac{|t|^N}{\sin\delta}
 \sum_{r\geq1}\rho^r r^{N-1}
 \cdot 2\zeta(N+1)\left(\frac{m}{2\pi}\right)^{N+1}.
\]
Here $\zeta(s)$ denotes the Riemann zeta function, not the root of unity.
Since
\[
 \sum_{r\geq1}e^{-\ell r}r^{N-1}
 \leq e^{\ell}\int_0^{\infty}e^{-\ell x}x^{N-1}\,dx
 =\frac{(N-1)!}{\rho\ell^N},
\]
and $\zeta(N+1)\leq\zeta(2)=\pi^2/6$, the stated constant follows.
The integral estimate can be checked termwise by integrating on
$[r,r+1]$; enlarging $[1,\infty)$ to $[0,\infty)$ only increases it.

For real $\xi$, all exponentials $e^{\xi/\omega_j}$ have modulus one.
Thus the absolute value of the rational factor in
\eqref{eq:cb-borel} is at most $\rho/(1-\rho^m)$, while
$\sum_{j\ne0}|\omega_j|^{-2}=m^2/12$.
For the horizontal strip, use
\[
 |a^m e^{m\xi/\omega_j}|
 \leq \rho^m\exp\!\left(\frac{m^2b}{2\pi|j|}\right)
 \leq \rho^m\exp\!\left(\frac{m^2b}{2\pi}\right)<1.
\]
The numerator is uniformly bounded as well, proving the strip assertion.

Finally, on the positive real Borel ray the double series before the
geometric summation is absolutely dominated by
$\sum_r\rho^r\sum_{j\ne0}|\omega_j|^{-2}$.
Fubini's theorem and $\Re(1/t)>0$ give
\[
 \int_0^{\infty}e^{-\xi/t}e^{r\xi/\omega_j}\,d\xi
 =\frac{t}{1-rt/\omega_j}.
\]
Equation \eqref{eq:cb-exact-kernel} now proves
\eqref{eq:cb-positive-laplace}.
\end{proof}

\subsection{Admissible rays and accumulation of Stokes directions}

For a ray avoiding \eqref{eq:cb-poles}, define
\begin{equation}
 \mathcal S_\theta\widetilde R_a(t)
 =\int_{0}^{e^{i\theta}\infty}
           e^{-\xi/t}B_{a,\zeta}(\xi)\,d\xi.
 \label{eq:cb-directional-sum}
\end{equation}
For each fixed $a$ and each such direction, the Borel transform is bounded
on the ray, so this integral converges whenever
$\Re(e^{i\theta}/t)>0$.
Indeed there are only finitely many profiles
\[
 F_p(w)=\frac{a^p e^{pw}}{1-a^m e^{mw}},
 \qquad p=1,\ldots,m,
\]
and two ray orientations for $w=\xi/\omega_j$.  If $\Re w\to-\infty$
the profile tends to zero.  If $\Re w\to+\infty$, it is asymptotic to
$-a^{p-m}e^{(p-m)w}$, which is bounded because $p\leq m$.  If
$\Re w=0$, the denominator is bounded below by $1-|a|^m$.  In the
remaining compact part of each profile ray, avoidance of its zeros gives
a positive minimum for the denominator.  Summing $O(|\omega_j|^{-2})$
proves the assertion.  No Diophantine condition is being tacitly assumed;
the bound is for fixed $a$ and a fixed nonsingular ray, and need not stay
bounded when the ray or parameter approaches a singular configuration.

This raywise statement requires a distinction from a theorem asserting
analyticity in an infinite angular Borel sector.  For fixed $a\ne0$, the
directions of \eqref{eq:cb-poles} accumulate at both real directions as
$n\to\pm\infty$.  In particular, although the positive real ray itself
is nonsingular and \eqref{eq:cb-positive-laplace} holds, there is no open
angular sector about that ray which is free of all Borel poles.  Its
horizontal-strip continuation is the relevant substitute.  The poles
are locally finite in the Borel plane; it is their \emph{directions} that
accumulate, at points tending to infinity.  These statements are
compatible, and they should not be conflated with the assumptions on a
finite collection of actions used in a particular transseries chart.

Whenever a deformation from an admissible contour $\gamma_-$ to
$\gamma_+$ sweeps the poles in a set $P$, with vanishing closing
integrals and an absolutely convergent residue sum, the residue theorem
gives the exact identity
\begin{equation}
 \mathcal S_+\widetilde R_a(t)-\mathcal S_-\widetilde R_a(t)
 =\sum_{(j,n):\,\xi_{j,n}\in P}
       \frac{e^{2\pi i n p_j/m}}{j}e^{-\xi_{j,n}/t}.
 \label{eq:cb-stokes-general}
\end{equation}
The orientation convention is that the contour formed by
$\gamma_-$ followed by $-\gamma_+$ encloses $P$ counterclockwise.
Equivalently, its right side is
$-2\pi i\sum_{\xi\in P}\operatorname{Res}_\xi
 [e^{-\xi/t}B_{a,\zeta}(\xi)]$.
The contour hypotheses are stated explicitly; the next proposition
checks them with elementary estimates in the case needed to disprove
the repository's assertion.

\subsection{A concrete counterexample to direction-independent normalization}

\begin{proposition}[Two nonsingular rays give different sums]
\label{thm:direction-counterexample}
Take $m=1$, $\zeta=1$, and $a=e^{-1}$.  Both the positive real ray and
the ray of angle $\pi/4$ avoid every Borel singularity, and their
Laplace integrals converge for every $t>0$.  They satisfy
\begin{align}
 \mathcal S_{\pi/4}\widetilde R_a(t)-R_a(t)
 &=\sum_{j=1}^{\infty}\sum_{n=1}^{\infty}
       \frac{1}{j}
       \exp\!\left[-\frac{j(4\pi^2n+2\pi i)}{t}\right]
       \label{eq:cb-counterexample-sum}\\
 &=-\sum_{n=1}^{\infty}
       \log\!\left(1-
          \exp\!\left[-\frac{4\pi^2n+2\pi i}{t}\right]\right).
       \label{eq:cb-counterexample-product}
\end{align}
The logarithms in the second line are their convergent power-series
branches.  The difference is nonzero for every $t=1/N$, $N\geq1$ an
integer, and for all sufficiently small positive $t$.
\end{proposition}

\begin{proof}
For $m=1$, pairing the summands $j$ and $-j$ in
\eqref{eq:cb-borel} gives
\begin{equation}
 B_{a,1}(\xi)=-\frac{1}{4\pi^2}
  \sum_{j=1}^{\infty}\frac{1}{j^2}
  \left(
   \frac{ae^{i\xi/(2\pi j)}}{1-ae^{i\xi/(2\pi j)}}
   +\frac{ae^{-i\xi/(2\pi j)}}{1-ae^{-i\xi/(2\pi j)}}
  \right).
 \label{eq:cb-mone-borel}
\end{equation}
The poles in the open wedge $0<\arg\xi<\pi/4$ are precisely
\[
 \xi=4\pi^2jn+2\pi i j,\qquad j,n\geq1,
\]
each with residue $i/(2\pi j)$.  No pole lies on either boundary ray:
in the upper half-plane a pole on the diagonal would require
$4\pi^2jn=2\pi j$, or $2\pi n=1$.

To justify contour rotation despite the infinitely many Stokes directions
near zero, rotate \emph{each paired $j$th summand} first.  Put
$\xi=2\pi j u$.  Close the sector from the positive real $u$-axis to
the diagonal along the vertical line
$\Re u=(2M+1)\pi$, with $M\to\infty$.  On this vertical segment both
numbers $ae^{iu}$ and $ae^{-iu}$ are negative real.  Each quotient
$ae^{\pm iu}/(1-ae^{\pm iu})$ therefore has modulus at most one.
For $t>0$, the modulus of the exponential kernel there is
$\exp[-2\pi j(2M+1)\pi/t]$; the segment length grows only linearly in
$M$.  The closing integral consequently tends to zero.

The residue theorem, with the indicated orientation, gives for this
paired summand
\[
 (\mathcal S_{\pi/4}-\mathcal S_0)_j
 =\frac{1}{j}\sum_{n\geq1}
          e^{-j(4\pi^2n+2\pi i)/t}.
\]
The integrals over either infinite ray can be summed over $j$ by the
bounded-profile estimate and $\sum j^{-2}<\infty$; the displayed
residue sums are absolutely convergent.  Summing proves
\eqref{eq:cb-counterexample-sum}, because
$\mathcal S_0\widetilde R_a=R_a$ by
\eqref{eq:cb-positive-laplace}.  Summation first over $j$ gives
\eqref{eq:cb-counterexample-product}.

For $t=1/N$ all factors $e^{-2\pi i j/t}$ equal one, so every term in
\eqref{eq:cb-counterexample-sum} is positive real.  More generally, as
$t\downarrow0$,
\[
 \mathcal S_{\pi/4}\widetilde R_a(t)-R_a(t)
 =e^{-(4\pi^2+2\pi i)/t}
       \bigl(1+O(e^{-4\pi^2/t})\bigr),
\]
since every other pair $(j,n)$ has $jn\geq2$.  This also proves
nonvanishing for all sufficiently small positive $t$.
\end{proof}

Consequently the conjecture's uniform Gevrey assertion, meromorphic
continuation assertion, and positive-ray reconstruction are valid in
the stated fixed-parameter disk.  Its demand that \emph{every}
nonsingular convergent directional sum equal the single fixed product
is refuted.  One can either keep the positive-ray normalization or attach
the explicit Stokes correction to each other contour.  Keeping only the
algebraic expansion cannot distinguish these functions.

\subsection{The corresponding correction for an inverse branch}

The preceding computation does not, by itself, establish every functional
analytic hypothesis in a parameter-dependent resurgent implicit-function
theorem.  It does establish the exact transformation that any such
theorem must respect.  Suppose on a common sector that
\[
 F_+(a,t)=F_-(a,t)+D(a,t),
\]
where $F_\pm$ are two sectorial determinations of $L_a$ and $D$ is the
residue correction.  Let $v_\pm(t)$ be the corresponding sectorial data,
and let $a_\pm(t)$ be chosen local solutions of
\[
 F_\pm(a_\pm,t)=F_\pm(a_0,t)+v_\pm(t).
\]
Set
\begin{align*}
 H_-(a,t)&=F_-(a,t)-F_-(a_0,t)-v_-(t),\\
 E(a,t)&=D(a,t)-D(a_0,t)-(v_+(t)-v_-(t)).
\end{align*}
Then the exact relation is
\begin{equation}
 H_-(a_+(t),t)=-E(a_+(t),t).
 \label{eq:cb-inverse-exact}
\end{equation}
If $A_-(w,t)$ denotes the chosen local inverse of $a\mapsto H_-(a,t)$,
this is equivalently
\begin{equation}
 a_+(t)=A_-\bigl(-E(a_+(t),t),t\bigr),
 \qquad a_-(t)=A_-(0,t).
 \label{eq:cb-inverse-map}
\end{equation}
Existence and uniqueness here are local analytic statements, requiring
nonvanishing of $\partial_aH_-$ and sufficiently small perturbations
within the chosen inverse chart.

For a precise first-order version, insert a scalar parameter $\epsilon$
and solve $H_-(a_\epsilon,t)+\epsilon E(a_\epsilon,t)=0$ near
$\epsilon=0$.  Differentiating the ordinary implicit equation gives
\begin{equation}
 \left.\frac{\partial a_\epsilon(t)}{\partial\epsilon}
 \right|_{\epsilon=0}
 =-\frac{E(a_-(t),t)}{\partial_aH_-(a_-(t),t)}.
 \label{eq:cb-inverse-first-jump}
\end{equation}
This formulation avoids an unproved uniform error estimate when either
derivatives or inverse radii depend on $t$.  It also makes the
normalization essential: the correction at the reference point $a_0$
and any jump in the target $v$ must both be included.

For inversion with a fixed target $y$ rather than a subtracted reference
value, the same statement reads
\[
 F_+(\,\cdot\,,t)^{-1}
   =\bigl(F_-(\,\cdot\,,t)+D(\,\cdot\,,t)\bigr)^{-1}
\]
on the chosen local branch.  Thus, whenever a resurgent algebra and its
summation maps satisfy the required composition and implicit-function
closure hypotheses, the Stokes map acts on the \emph{equation and its
target}, and the transformed equation is then inverted.  Inversion does
not remove the Stokes correction.

